\documentclass[17pt, aspectratio=169]{beamer}

% --- Paquetes de Idioma y Estética ---
\usepackage[utf8]{inputenc}
\usepackage[T1]{fontenc}
\usepackage[spanish]{babel}
\usepackage{xcolor}
\usepackage{tcolorbox}
\tcbuselibrary{skins}

% --- Definición de Colores ICMG ---
\definecolor{ICMGBlue}{HTML}{002366}
\definecolor{ICMGGold}{HTML}{996515}
\definecolor{CitationBg}{HTML}{F4F4F4}

% --- Configuración del Tema Beamer ---
\setbeamercolor{palette primary}{bg=ICMGBlue, fg=white}
\setbeamercolor{palette secondary}{bg=ICMGBlue, fg=white}
\setbeamercolor{structure}{fg=ICMGBlue}
\setbeamercolor{title}{fg=ICMGBlue}
\setbeamerfont{title}{series=\bfseries}

\setbeamertemplate{navigation symbols}{}

\setbeamertemplate{footline}{
  \leavevmode%
  \hbox{%
  \begin{beamercolorbox}[wd=.5\paperwidth,ht=2.25ex,dp=1ex,left]{author in head/foot}%
    \hspace*{2ex} Clase 1: El Evangelio
  \end{beamercolorbox}%
  \begin{beamercolorbox}[wd=.5\paperwidth,ht=2.25ex,dp=1ex,right]{date in head/foot}%
    Iglesia Cristiana Manantial de la Gracia \hspace*{2ex} \insertframenumber{} / \inserttotalframenumber \hspace*{2ex} 
  \end{beamercolorbox}}%
  \vspace{2pt}
}

% --- Formato para Citas Bíblicas ---
\newtcolorbox{citaBiblicaSlide}[1]{
  enhanced,
  colback=CitationBg,
  colframe=ICMGGold,
  leftrule=3pt,
  rightrule=0pt,
  toprule=0pt,
  bottomrule=0pt,
  arc=0mm,
  fonttitle=\bfseries\color{ICMGBlue},
  title={#1},
  boxrule=0pt,
  italic
}

\title{Clase 1: El Evangelio}
%\subtitle{Clase 1: El Fundamento de Nuestra Identidad}
\author{Iglesia Cristiana Manantial de la Gracia}
\date{Serie: El Bautismo}

% --- Configuración de Diapositiva de Título Intermedia ---
\AtBeginSection[]{
  \begin{frame}[plain] % 'plain' quita el pie de página para resaltar el título
    \centering
    \vfill
    {\color{ICMGGold}\rule{0.6\textwidth}{1.5pt}} % Línea decorativa superior
    \vspace{0.4cm}
    
    {\color{ICMGBlue}\scshape\large SECCIÓN \thesection \par} % "SECCIÓN 1", etc.
    \vspace{0.2cm}
    {\color{ICMGBlue}\Huge\bfseries \insertsectionhead \par} % Nombre de la sección
    
    \vspace{0.4cm}
    {\color{ICMGGold}\rule{0.6\textwidth}{1.5pt}} % Línea decorativa inferior
    \vfill
  \end{frame}
}

\begin{document}

\begin{frame}[plain]
    \titlepage
    \centering
    \color{ICMGGold}\rule{0.5\textwidth}{1pt}
\end{frame}

\section{¿Qué es el bautismo?}

\begin{frame}{¿Qué es el bautismo?}
    \begin{itemize}
        \item Medio de Gracia (no salva).
    \end{itemize}
\end{frame}

\begin{frame}{¿Qué es el bautismo?}
    \begin{itemize}
        \item Medio de Gracia (no salva).
        \item Ordenanza (Gran Comisión)
    \end{itemize}
\end{frame}

\begin{frame}{¿Qué es el bautismo?}
    \begin{itemize}
        \item Medio de Gracia (no salva).
        \item Ordenanza (Gran Comisión).
        \item Testimonio público de alguien que profesa fe en Cristo.
    \end{itemize}
\end{frame}

%\begin{frame}{¿Qué es el bautismo?}
    %\begin{itemize}
        %\item Medio de Gracia (no salva).
        %\item Reconocer que la Biblia es la única %autoridad y estándar infalible.
    %\end{itemize}
    
    %\begin{citaBiblicaSlide}{Juan 17:3}
        %"Y esta es la vida eterna: que te conozcan a ti, %el único Dios verdadero, y a Jesucristo, a quien %has enviado."
    %\end{citaBiblicaSlide}
    
%\end{frame}
%Una diapositiva hablando de la Gran Comisión: Mandato de Jesús por lo que nosotros también lo tenemos que cumplir y que, cuando nos enseñan lo que Él ha mandado, nosotros debemos ser bautizados

%El resto de las diapositivas con un análisis del evangelio en Hechos 17: 22-31

\section{¿Qué involucra la Gran Comisi\'on?}

\begin{frame}{¿Qué involucra la Gran Comisi\'on?}
    \begin{itemize}
        \item Nosotros fuimos y somos discípulos de alguien.
    \end{itemize}
    
    \begin{citaBiblicaSlide}{Mateo 28:19}
        " Vayan, pues, y hagan \textbf{discípulos} de todas las naciones, \textbf{bautizándolos} en el nombre del Padre y del Hijo y del Espíritu Santo,"
    \end{citaBiblicaSlide}
\end{frame}

\begin{frame}{¿Qué involucra la Gran Comisi\'on?}
    \begin{itemize}
        \item El ser disc\'ipulo implica ser bautizado y enseñado (continuamente) (cf. Hechos 2:38).
    \end{itemize}
    
    \begin{citaBiblicaSlide}{Mateo 28:20}
        "\textbf{enseñándoles} a guardar todo lo que les he mandado; y ¡recuerden! Yo estoy con ustedes todos los días, hasta el fin del mundo."
    \end{citaBiblicaSlide}
\end{frame}

\section{El Evangelio: Hechos 17:24-31}
%Aquí lo mejor creo que es ponerlo en diapositivas normales, no en doble columna y rescatando lo que ya tengo acá abajo, sólo que en mejor formato

%Agregar conclusión chida de lo que me escupió Logos, algo que no es sólo de conocimeinto teórico: La implicación práctica es que guardar los mandamientos de Jesús no es una actividad pasiva o meramente intelectual, sino un proceso dinámico de transformación comunitaria que modela la compasión, la justicia y la reconciliación que Jesús encarnó y enseñó.

\begin{frame}{El Evangelio: Hechos 17:24-31}
    \begin{itemize}
        \item Dios como creador soberano.
    \end{itemize}
    
    \begin{citaBiblicaSlide}{Hechos 17:24}
        "»El Dios que hizo el mundo y todo lo que en él hay, puesto que es Señor del cielo y de la tierra, no mora en templos hechos por manos de hombres,"
    \end{citaBiblicaSlide}
\end{frame}

\begin{frame}{El Evangelio: Hechos 17:24-31}
    \begin{itemize}
        \item Dios como creador soberano.
        \item Dios como sustentador de la vida.
    \end{itemize}
    
    \begin{citaBiblicaSlide}{Hechos 17:25}
        "ni es servido por manos humanas, como si necesitara de algo, puesto que Él da a todos vida y aliento y todas las cosas."
    \end{citaBiblicaSlide}
\end{frame}

\begin{frame}{El Evangelio: Hechos 17:24-31}
    \begin{itemize}
        \item Dios como creador soberano.
        \item Dios como sustentador de la vida.
        \item \textbf{El} propósito del hombre (v.26-28).
    \end{itemize}
\end{frame}

\begin{frame}{El Evangelio: Hechos 17:24-31}
    \begin{itemize}
        \item Dios como creador soberano.
        \item Dios como sustentador de la vida.
        \item \textbf{El} propósito del hombre (v.26-28).
        \item La brecha entre Dios y el hombre (v.29-30).
    \end{itemize}
\end{frame}

\begin{frame}{El Evangelio: Hechos 17:24-31}
    \begin{itemize}
        \item Dios como creador soberano.
        \item Dios como sustentador de la vida.
        \item \textbf{El} propósito del hombre (v.26-28).
        \item La brecha entre Dios y el hombre (v.29-30).
        \item La justicia de Dios (v.31).
    \end{itemize}
\end{frame}

\begin{frame}{Conclusi\'on}
    \begin{itemize}
        \item El conocimiento no salva.
    \end{itemize}
\end{frame}

\begin{frame}{Conclusi\'on}
    \begin{itemize}
        \item El conocimiento no salva.
        \item El verdadero arrepintimiento trae fruto (Efesios 4:1-6).
    \end{itemize}
\end{frame}

\end{document}